\begin{definition} \label{definition:self_adjoint_hermitian_skew_symmetric_unitary_normal_operators}
    Let $V$ be a finite-dimensional inner product space over $\mathbb{R}$ or $\mathbb{C}$, and let $T: V \to V$ be a linear endomorphism.

    \begin{itemize}
        \item $T$ is called \hldef{self-adjoint} (or \hldef{Hermitian}, in the complex case) if $T^* = T$.
        \item $T$ is called \hldef{skew-adjoint} (or \hldef{skew-Hermitian}) if $T^* = -T$.
        \item $T$ is called \hldef{unitary} (resp. \hldef{orthogonal}, in the real case) if $T^* T = T T^* = I_V$.
        \item $T$ is called \hldef{normal} if $T^* T = T T^*$.
    \end{itemize}

    In terms of an orthonormal basis, these conditions correspond to the matrix being equal to its conjugate transpose, negative conjugate transpose, having conjugate transpose as inverse, or commuting with its conjugate transpose, respectively.
\TextIfExists{definition:involution_on_a_ring}{These classifications of operators are special cases of elements in a ring with \Cref{definition:involution_on_a_ring}: self-adjoint ($x^* = x$), skew-adjoint ($x^* = -x$), unitary ($x^*x = xx^* = 1$), and normal ($x^*x = xx^*$).}
\TODO{AI generated, verify}
\end{definition}
